\honest{0 And 1 Are Not Probabilities}{Eliezer Yudkowsky}{2008}

Infinity is not an integer. You never reach it by counting, and it does not obey the rules that integers obey: ``5 + infinity = infinity'' does not give you ``infinity - infinity = 5.''

Probabilities are usually written between 0 and 1, but they can be rewritten as odds, $P/(1-P)$, and back again. Odds make Bayesian updates easy. If a bell is twice as likely after a roll of 1 as after other rolls of a die, hearing it moves the odds of a 1 from 1:5 to 2:5, a probability of 2/7. Cox's theorem, I say, shows that all reasonable ways of representing uncertainty are equivalent. \nb{In Jaynes's presentation of that theorem, ``Certainty is represented by w(A|C) = 1.'' The theorem puts certainty on the scale.}

On odds, 1 becomes infinity, so ``absolute truth doesn't look like it should be so easy to reach.'' \nb{On odds, 0 still goes to 0, so impossibility looks as easy to reach as before. What looks reachable depends on the scale.} On log odds, measured in decibels as Jaynes recommended, 0 and 1 are both infinite. There the distance between two degrees of belief equals the evidence needed to move from one to the other, and ``reaching infinite certainty requires infinitely strong evidence.'' \nb{This is correct, and it is the useful point of the post.}

Theorems of probability also need special cases at 0 and 1, such as updating on an observation to which you assigned probability 0. \nb{Standard theory meets that case constantly. For a quantity with a continuous distribution, every exact value has probability 0, yet one of them occurs, and updating on it is done with densities. In standard theory, probability 0 does not mean impossible.}

So I propose that 0 and 1 are not probabilities, just as infinity is not a real number. \nb{The argument showed that certainty cannot be reached by updating, not that nothing may have probability 1. Taken literally, the proposal would also forbid continuous distributions, since at most countably many exclusive events can have positive probability.} Probability theorists would object, I admit, because theorems that assume the probabilities of all possibilities add up to 1 would have to be rederived. \nb{That assumption is one of Kolmogorov's axioms, and it gives probability 1 to $A$ or not-$A$ in every model. So the analogy with infinity does not hold: infinity breaks the axioms of the real numbers, while the axioms of probability require 0 and 1.}

But a real die can land on its edge. If you invent a ``magical symbol'' for all the possibilities you have not considered, you arrive at a symbol for infinite certainty. \nb{Standard theory calls this the sample space. It is an ordinary event with probability 1.} I would rather find a way to derive the theorems without magic symbols. Some mathematicians refuse to believe in the law of the excluded middle, and I would like to be a probability theorist who does not believe in absolute certainty. I do not say how.
